\section*{Приложение A. Воспроизводимость экспериментов}
\addcontentsline{toc}{section}{Приложение A. Воспроизводимость экспериментов}
\setcounter{table}{0}
\setcounter{figure}{0}
\setcounter{equation}{0}
\renewcommand{\thetable}{A.\arabic{table}}
\renewcommand{\thefigure}{A.\arabic{figure}}
\renewcommand{\theequation}{A.\arabic{equation}}

Данное приложение фиксирует технические параметры, использованные файлы, артефакты и команды воспроизводимости для трёх экспериментальных блоков: основного эксперимента RUN 03, дополнительного эксперимента 1 RUN 04 и дополнительного эксперимента 2 RUN 05. Цель приложения --- сделать экспериментальную часть проверяемой: читатель должен понимать, какие данные, признаки, модели и результаты использовались при формировании выводов курсовой работы.

\subsection*{Структура проекта}

Код и материалы экспериментов хранятся в репозитории:

\begin{Verbatim}[fontsize=\small]
MaksimDugin/forecast-with-llm-features
\end{Verbatim}

Ключевые элементы проекта приведены в таблице~\ref{tab:appendix_project_structure}.

\begin{table}[htbp]
\centering
\begin{tabular}{ll}
\toprule
Путь & Назначение \\
\midrule
\texttt{notebooks/} & Jupyter notebooks для подготовки данных и запуска экспериментов \\
\texttt{experiments/} & Python scripts для воспроизводимых запусков \\
\texttt{artifacts/} & выходные таблицы, метрики, предсказания и служебные файлы \\
\texttt{reports/} & Markdown-отчёты по данным и результатам \\
\texttt{docs/} & методологические заметки \\
\texttt{coursework\_latex/} & LaTeX-версия текста курсовой \\
\texttt{coursework\_latex/sections/} & отдельные LaTeX-разделы курсовой \\
\bottomrule
\end{tabular}
\caption{Основная структура проекта}
\label{tab:appendix_project_structure}
\end{table}

Для текста курсовой используются следующие LaTeX-файлы:

\begin{Verbatim}[fontsize=\small]
coursework_latex/main_msu_template.tex
coursework_latex/sections/00_introduction.tex
coursework_latex/sections/01_literature_review.tex
coursework_latex/sections/02_experiment_setup.tex
coursework_latex/sections/02_3_features.tex
coursework_latex/sections/02_4_methodology.tex
coursework_latex/sections/02_5_models.tex
coursework_latex/sections/02_6_results.tex
coursework_latex/sections/02_8_run04_additional_experiment_template.tex
coursework_latex/sections/02_9_run05_multi_asset_experiment.tex
coursework_latex/sections/02_7_interpretation_limitations.tex
coursework_latex/sections/03_conclusion.tex
coursework_latex/sections/04_appendix_reproducibility.tex
coursework_latex/references_draft.bib
\end{Verbatim}

\subsection*{Общая логика воспроизводимости}

Все эксперименты строятся как проверка групп признаков, а не как соревнование архитектур. Поэтому при воспроизведении важно сохранить не только итоговые модели, но и одинаковый протокол подготовки данных:

\begin{enumerate}
    \item целевая переменная строится только из будущей цены закрытия и текущей цены закрытия;
    \item данные сортируются по дате принятия решения;
    \item обучающая, валидационная и тестовая подвыборки не пересекаются по датам;
    \item текстовые признаки строятся только из информации, доступной к моменту прогноза;
    \item scaler, PCA и другие преобразования обучаются только на обучающей подвыборке;
    \item гиперпараметры и порог классификации выбираются только на валидационной подвыборке;
    \item тестовая подвыборка используется только для финальной оценки;
    \item результаты сравниваются не только по accuracy, но и по balanced accuracy, F1 и ROC-AUC;
    \item вывод формулируется через влияние групп признаков, а не через качество одной модели.
\end{enumerate}

\subsection*{RUN 03: основной эксперимент AAPL}

Основной эксперимент RUN 03 используется как первичная экспериментальная проверка. Он сравнивает три укрупнённые группы признаков: \texttt{prices-only}, \texttt{FinBERT-only} и \texttt{prices + FinBERT}. Параметры запуска приведены в таблице~\ref{tab:appendix_run03_params}.

\begin{table}[htbp]
\centering
\begin{tabular}{ll}
\toprule
Параметр & Значение \\
\midrule
Актив & AAPL \\
Частота данных & дневная \\
Горизонт прогноза & 1 торговый день \\
Тип задачи & бинарная классификация \\
Целевая переменная & $y_t = \mathbf{1}\{Close_{t+1} > Close_t\}$ \\
Текстовый лаг & $text \leq t-1$ \\
Модель embeddings & \texttt{ProsusAI/finbert} \\
Размерность embedding & 768 \\
Число рыночных признаков & 11 \\
Число наблюдений & 1578 \\
Период данных & 2020-04-01 --- 2020-09-30 \\
Разбиение & \texttt{chronological\_row\_balanced\_by\_decision\_date} \\
\bottomrule
\end{tabular}
\caption{Параметры основного эксперимента RUN 03}
\label{tab:appendix_run03_params}
\end{table}

Целевая переменная задаётся формулой:

\[
y_t = \mathbf{1}\{Close_{t+1} > Close_t\}.
\]

Рыночный вектор признаков имеет вид

\[
p_t \in \mathbb{R}^{11},
\]

FinBERT embedding имеет вид

\[
e_t = FinBERT(text_t) \in \mathbb{R}^{768},
\]

а объединённый вектор признаков задаётся как

\[
x_t = [p_t;e_t] \in \mathbb{R}^{779}.
\]

В RUN 03 используется хронологическое разбиение по дате принятия решения. Обучающая, валидационная и тестовая подвыборки не пересекаются по датам.

\begin{table}[htbp]
\centering
\begin{tabular}{lrrl}
\toprule
Подвыборка & Наблюдения & Уникальные даты & Период \\
\midrule
Train & 1104 & 52 & 2020-04-01 --- 2020-06-15 \\
Validation & 252 & 10 & 2020-06-16 --- 2020-06-29 \\
Test & 222 & 55 & 2020-06-30 --- 2020-09-30 \\
\bottomrule
\end{tabular}
\caption{Разбиение данных в RUN 03}
\label{tab:appendix_run03_split}
\end{table}

\subsection*{RUN 04: расширенный рыночный baseline и PCA}

Дополнительный эксперимент 1 RUN 04 проверяет, сохраняется ли вклад FinBERT embeddings после усиления рыночной базовой модели. Он использует расширенные рыночные признаки, rolling-признаки, volume-based признаки и PCA-сжатые FinBERT embeddings.

\begin{table}[htbp]
\centering
\begin{tabular}{ll}
\toprule
Параметр & Значение \\
\midrule
Актив & AAPL \\
Источник данных & \texttt{TheFinAI/flare-sm-bigdata} \\
Частота данных & дневная \\
Горизонт прогноза & 1 торговый день \\
Число текстов & 497 \\
Число наблюдений & 163 \\
Число уникальных дат & 163 \\
Период данных & 2020-03-09 --- 2020-10-27 \\
Текстовое окно & до 10 дней \\
Текстовый лаг & тексты не позднее $t-1$ \\
Размерность FinBERT embedding & 768 \\
PCA-размерности & $k \in \{16,32,64\}$ \\
\bottomrule
\end{tabular}
\caption{Параметры дополнительного эксперимента RUN 04}
\label{tab:appendix_run04_params}
\end{table}

Хронологическое разбиение RUN 04 имеет вид:

\begin{table}[htbp]
\centering
\begin{tabular}{lrrl}
\toprule
Подвыборка & Наблюдения & Даты & Период \\
\midrule
Train & 114 & 114 & 2020-03-09 --- 2020-08-18 \\
Validation & 24 & 24 & 2020-08-19 --- 2020-09-22 \\
Test & 25 & 25 & 2020-09-23 --- 2020-10-27 \\
\bottomrule
\end{tabular}
\caption{Разбиение данных в RUN 04}
\label{tab:appendix_run04_split}
\end{table}

Ключевой script для RUN 04:

\begin{Verbatim}[fontsize=\small]
experiments/run_04_feature_ablation.py
\end{Verbatim}

Основные выходные артефакты RUN 04:

\begin{table}[htbp]
\centering
\begin{tabular}{ll}
\toprule
Файл & Содержание \\
\midrule
\texttt{artifacts/run\_04\_load\_dataset/metrics.csv} & итоговые test-метрики по группам признаков \\
\texttt{artifacts/run\_04\_load\_dataset/validation\_selection.csv} & выбор $C$, порога и конфигурации на validation \\
\texttt{artifacts/run\_04\_load\_dataset/dataset\_summary.json} & сводка по данным, датам, PCA и размерам групп признаков \\
\bottomrule
\end{tabular}
\caption{Основные артефакты RUN 04}
\label{tab:appendix_run04_artifacts}
\end{table}

Пример запуска script для готового CSV-файла:

\begin{Verbatim}[fontsize=\small]
python experiments/run_04_feature_ablation.py \
  --input path/to/ready_dataset.csv \
  --output-dir artifacts/run_04_load_dataset \
  --date-col date \
  --target-col y \
  --ticker-col asset \
  --embedding-prefix finbert_
\end{Verbatim}

\subsection*{RUN 05: мультиактивная проверка crypto-news}

Дополнительный эксперимент 2 RUN 05 проверяет, сохраняются ли выводы при переходе к другому источнику текстов, другому классу активов и нескольким горизонтам прогноза. В нём используются новости \texttt{crypto-news} и дневные OHLCV-данные Binance.

\begin{table}[htbp]
\centering
\begin{tabular}{ll}
\toprule
Параметр & Значение \\
\midrule
Источник текстов & \texttt{crypto-news} \\
Источник цен & Binance daily OHLCV \\
Активы & BNB, BTC, SOL \\
Горизонты & 1, 3, 7 календарных дней \\
Текстовые признаки & sentiment-score, sentiment probabilities, text-meta \\
Рыночная группа & \texttt{market\_full} \\
Характер эксперимента & exploratory robustness check \\
\bottomrule
\end{tabular}
\caption{Параметры дополнительного эксперимента RUN 05}
\label{tab:appendix_run05_params}
\end{table}

Целевая переменная для RUN 05 задаётся как

\[
y_{t,h}=\mathbf{1}\{Close_{t+h}>Close_t\}, \qquad h \in \{1,3,7\}.
\]

Покрытие итоговых датасетов после объединения новостей с ценами:

\begin{table}[htbp]
\centering
\begin{tabular}{lrr}
\toprule
Актив & Наблюдения & Уникальные даты \\
\midrule
BNB & 1873 & 569 \\
BTC & 6000 & 752 \\
SOL & 885 & 476 \\
\bottomrule
\end{tabular}
\caption{Покрытие данных в RUN 05}
\label{tab:appendix_run05_coverage}
\end{table}

Ключевой notebook для RUN 05:

\begin{Verbatim}[fontsize=\small]
run05_multi_dataset_pipeline.ipynb
\end{Verbatim}

Он выполняет следующие шаги:

\begin{enumerate}
    \item загружает \texttt{crypto-news};
    \item фильтрует новости по активам BNB, BTC и SOL;
    \item загружает дневные OHLCV-данные Binance;
    \item строит целевую переменную для горизонтов 1, 3 и 7 дней;
    \item извлекает FinBERT-derived sentiment-признаки и text-meta признаки;
    \item формирует готовые CSV-файлы для оценки;
    \item запускает единый evaluation engine на основе \texttt{experiments/run\_04\_feature\_ablation.py};
    \item сохраняет итоговые метрики и дельты относительно \texttt{market\_full}.
\end{enumerate}

Основные выходные артефакты RUN 05:

\begin{table}[htbp]
\centering
\begin{tabular}{ll}
\toprule
Файл & Содержание \\
\midrule
\texttt{artifacts/run\_05\_multi\_dataset\_pipeline/run05\_run\_log.csv} & статус запуска по активам и горизонтам \\
\texttt{artifacts/run\_05\_multi\_dataset\_pipeline/run05\_all\_metrics.csv} & все test-метрики по всем конфигурациям \\
\texttt{artifacts/run\_05\_multi\_dataset\_pipeline/run05\_text\_deltas.csv} & приросты текстовых признаков относительно \texttt{market\_full} \\
\texttt{artifacts/run\_05\_multi\_dataset\_pipeline/ready\_datasets/} & готовые CSV-файлы для каждого актива и горизонта \\
\texttt{artifacts/run\_05\_multi\_dataset\_pipeline/run04\_results/} & результаты evaluation engine по каждому запуску \\
\bottomrule
\end{tabular}
\caption{Основные артефакты RUN 05}
\label{tab:appendix_run05_artifacts}
\end{table}

\subsection*{Команды воспроизводимости}

Минимальные зависимости для воспроизведения экспериментов:

\begin{Verbatim}[fontsize=\small]
pip install pandas numpy scikit-learn torch transformers tqdm tabulate kagglehub
\end{Verbatim}

Для запуска RUN 04 на готовом датасете используется script:

\begin{Verbatim}[fontsize=\small]
python experiments/run_04_feature_ablation.py \
  --input path/to/ready_dataset.csv \
  --output-dir artifacts/run_04_load_dataset \
  --date-col date \
  --target-col y \
  --ticker-col asset \
  --embedding-prefix finbert_
\end{Verbatim}

Для RUN 05 основной запуск выполняется из notebook:

\begin{Verbatim}[fontsize=\small]
run05_multi_dataset_pipeline.ipynb
\end{Verbatim}

Перед запуском RUN 05 рекомендуется задать отдельную папку для Kaggle cache, чтобы избежать переполнения системного диска:

\begin{Verbatim}[fontsize=\small]
from pathlib import Path
import os

KAGGLE_CACHE_DIR = r"D:\kagglehub_cache"
Path(KAGGLE_CACHE_DIR).mkdir(parents=True, exist_ok=True)
os.environ["KAGGLEHUB_CACHE"] = KAGGLE_CACHE_DIR
\end{Verbatim}

Для сборки LaTeX-документа используется корневой файл:

\begin{Verbatim}[fontsize=\small]
coursework_latex/main_msu_template.tex
\end{Verbatim}

Команды для локальной сборки:

\begin{Verbatim}[fontsize=\small]
cd coursework_latex
pdflatex main_msu_template.tex
bibtex main_msu_template
pdflatex main_msu_template.tex
pdflatex main_msu_template.tex
\end{Verbatim}

\subsection*{Краткий вывод по воспроизводимости}

Воспроизводимая часть работы фиксирует не только итоговые значения метрик, но и способ построения признаков, целевой переменной, хронологического разбиения и предварительной обработки. Это важно, поскольку в задачах финансового прогнозирования небольшие нарушения временного порядка или обучение preprocessing на будущих данных могут привести к завышенной оценке качества. Поэтому приложение является частью методологического контроля: оно показывает, какие именно условия должны быть выполнены, чтобы результаты можно было интерпретировать как проверку добавочной предсказательной силы LLM/FinBERT-generated features.
